% ============================================================
% §5  Conclusion
% ============================================================

\section{Conclusion}
\label{sec:conclusion}

This paper addresses the contaminated semi-supervised setting, in which training data contain a
small fraction of labeled anomalies alongside a majority of unlabeled observations.
This configuration is common in industrial deployments yet unaddressed by standard MTSAD
benchmarks or existing unsupervised methods.
CSMAD integrates labeled anomaly information into masked autoencoder representation learning
through co-designed, interacting paths --- loss bifurcation that restricts Student mimicry to
normal patches and gradient-reversal suppression of anomaly-specific information, with
anomaly-priority masking surfacing the labeled positions they act on. Together with the asymmetric
Teacher--Student decoder (3-layer Teacher, 2-layer Student), these paths turn the capacity gap into
a reliable, amplified discrepancy signal under contaminated training.
A contaminated benchmark protocol supports evaluation by incorporating the chronological prefix
of each test stream into training.
Experiments on [N] multivariate datasets demonstrate competitive performance against [N]
unsupervised and weakly supervised baselines under five complementary metrics; the label
sparsity analysis confirms graceful degradation as the labeled fraction decreases.
% PH:NUM-029 | dataset-family count (sync with NUM-001)
% PH:NUM-030 | total baseline count (sync with NUM-002)
A notable limitation is the cost of leave-one-out inference, which requires approximately
$50{\times}$ more forward-pass computation than single-mask scoring; reducing this cost, for
example through amortized inference with learned masking schedules or sparse patch selection,
is a direction for future work.
The graceful degradation toward the unsupervised limit also motivates a fully unsupervised
variant of CSMAD obtained by disabling the gradient-reversal pathway.
Code is available at [URL] (to be released upon acceptance).
% PH:TXT-002 | code repository URL
